% !PRINTED 545
\weberpage{545}

We now introduce the notation:
\weq{9}{\begin{aligned}
\eta &= \alpha + \alpha_e + \alpha_{2e} + \cdots + \alpha_{(f-1)e}\\
\eta_1 &= \alpha_1 + \alpha_{e+1} + \alpha_{2e+1} + \cdots + \alpha_{(f-1)e+1},\\
&\cdots\cdots\cdots\cdots\cdots\cdots\cdots\cdots\cdots\cdots\\
\eta_{e-1} &= \alpha_{e-1} + \alpha_{2e-1} + \alpha_{3e-1} + \cdots + \alpha_{m-1},
\end{aligned}}

and call these quantities, as \gesperrt{Gauss} did in the special case of the division of the circle, the $f$-term periods of the quantities $\alpha$. We then have
\weq{10}{(\varepsilon,\alpha) = \eta + \varepsilon\eta_1 + \varepsilon^2\eta_2 + \cdots + \varepsilon^{e-1}\eta_{e-1},}
for which we may also write $(\varepsilon,\eta)$.

The application of the permutation $\pi$ to the indices of $\alpha$ brings forth, among the indices of the quantities $\eta$, which are to be taken with respect to the modulus $e$, the cyclic permutation
\weq{11}{\gamma = (0, 1, 2\dots e-1)}

Corresponding to formula (10), we can now also write the formulas (5) and (7) thus:
\[
(\varepsilon,\eta)^{\nu} = E_0^{(\nu)} + \varepsilon E_1^{(\nu)} + \varepsilon^2 E_2^{(\nu)} + \cdots + \varepsilon^{e-1}E_{e-1}^{(\nu)}
\]
\[
(\varepsilon,\eta)^{\nu}(\varepsilon^{\lambda_1},\eta)^{\nu_1}(\varepsilon^{\lambda_2},\eta)^{\nu_2}\dots = G_0 + \varepsilon G_1 + \varepsilon^2 G_2 + \cdots + \varepsilon^{e-1}G_{e-1},
\]
in which then
\[
E_k^{(\nu)} = A_k^{(\nu)} + A_{e+k}^{(\nu)} + \cdots + A_{(f-1)e+k}^{(\nu)}
\]
\[
G_k = B_k + B_{e+k} + \cdots + B_{(f-1)e+k}
\]
and $E_k^{(\nu)}$ and $G_k$ are homogeneous polynomials in the $\alpha$, which can also be represented as functions of the $\eta$.

The quantities $\eta_k$, $E_k^{(\nu)}$, $G_k$ remain unchanged when the permutation $\pi^e$ is applied to their indices, that is, when the index is altered by a multiple of $e$; and we can therefore now complete Theorems 2. and 3. as follows:

\gesperrt{4. If $\varepsilon$ is an arbitrary $e$th root of unity and $e$ a divisor of $m$, then, when the permutation $\pi$ is applied to the indices of $\alpha$, the indices $k$ of the coefficients $E_k^{(\nu)}$ of $(\varepsilon,\eta)^{\nu}$ undergo the permutation $\gamma^{\nu}$, and the indices of the coefficients $G$ of the product $(\varepsilon,\eta)^{\nu}(\varepsilon^{\lambda_1},\eta)^{\nu_1}(\varepsilon^{\lambda_2},\eta)^{\nu_2}\dots$ the permutation $\gamma^{\nu+\lambda_1\nu_1+\lambda_2\nu_2\dots}$}
